\section{A uniform stationary chart and its auxiliary-volume derivatives}
\label{sec:uniform-stationary-chart}
We construct the stationary branch on the real coordinate neighborhood
used by the nonlinear step. The neighborhood must be uniform in the
auxiliary torus; the subsequent derivative bounds may pay a fixed
polynomial in its volume. These are different requirements.

\subsection{Local derivative data and the Hessian gap}
Let $I$ and $J$ be the coarse and free fine real-coordinate sets,
and write $w=(x,y)\in\mathbb R^I\times\mathbb R^J$. Both coordinate
spaces initially carry the supremum norm. Give every coordinate an
anchor $a_i$ on the fine torus, with at most $K$ coordinates at each
site. Put
\[
       S_a=K\sum_{r\ge0}(2r+1)^4e^{-ar}\qquad(a>0).
\]
The distance-shell bound implies
$\sum_i e^{-a\operatorname{dist}(v,a_i)}\le S_a$, uniformly in the
torus and in $v$.

Fix an order $N\ge3$. The local action estimates provide a smooth
real function $f(x,y)$ on $\|w\|_\infty<r_0$ and constants
$a,C,m,r_0>0$ such that $f(0)=0$, $D_yf(0)=0$, and
\begin{equation}\label{eq:stationary-scatter-input}
 |\partial_{i_1}\cdots\partial_{i_q}f(w)|
 \le C\exp\!\left(-a\sum_{h=1}^q
                  \operatorname{dist}(a_{i_t},a_{i_h})\right)
 \quad(1\le q\le N, 1\le t\le q).
\end{equation}
The vertical Hessian $H(0)=D_y^2f(0)$ satisfies
$H(0)\ge mI$ in the Euclidean quadratic-form sense.
The constants and coordinate multiplicity are independent of the
auxiliary torus after the block data have been fixed. The Hessian
gap is the tree-gauge energy comparison in
Section~\ref{subsec:path-stationary-chart}; the derivative estimate
is the additional quantitative local-action input. Its application here
uses the larger one-step data of that comparison: the coarse side exceeds
one and the fine side is at least $100R_{\rm macro}$, where
$R_{\rm macro}=(j_{\rm ph}+7)\ell$, the blocking side $\ell\ge3$ is odd,
the phase $j_{\rm ph}\ge4$, and the kernel's branch-compatibility
hypothesis holds.

For two points $v,w$ in the original coordinate ball, differentiate
a Hessian entry along their joining segment. Applying
\eqref{eq:stationary-scatter-input} with the first Hessian coordinate
as root, and summing the varied coordinate, gives
\[
 |H(v)_{ij}-H(w)_{ij}|
   \le C S_a\|v-w\|_\infty
                     e^{-a\operatorname{dist}(a_i,a_j)}.
\]
Both row and column sums of the difference are therefore at most
$CS_a^2\|v-w\|_\infty$. Its Euclidean operator norm and its
supremum operator norm have this same bound. Choose
\[
 r_H=\min\left\{r_0/2,\frac{m}{2(1+CS_a^2)}\right\}.
\]
Then $H(w)\ge(m/2)I$ on $\|w\|_\infty\le r_H$.
The second-derivative case of
\eqref{eq:stationary-scatter-input} also gives uniformly bounded
exponentially weighted row and column sums of $H(w)$, after
reducing $a$ to $a/2$. The inverse argument of
Section~\ref{sec:local-covariance}, applied with lower bound
$\min(1,m/2)$, supplies an exponential bound on $H(w)^{-1}$.
Summing that bound gives a fixed constant $J_0\ge1$ with
\begin{equation}\label{eq:stationary-uniform-inverse}
                 \|H(w)^{-1}\|_{\infty\to\infty}\le J_0.
\end{equation}
The Euclidean bound $\|H(w)^{-1}\|_{2\to2}\le2/m$ follows directly
from the positive gap. No conversion between the two vector norms
is used to obtain either uniform bound.

\subsection{Contraction on a common coordinate ball}
Let $F(x,y)=D_yf(x,y)$, viewed as its array of real coordinates.
Choose
\begin{align*}
 r_V&=\min\{r_H,[4J_0(1+CS_a^2)]^{-1}\},\\
 r_X&=\min\{r_V,r_V/[4J_0(1+CS_a)]\}.
\end{align*}
For $\|x\|_\infty<r_X$, define on the closed ball
$\|y\|_\infty\le r_V$
\[
                 T_x(y)=y-H(0)^{-1}F(x,y).
\]
The Hessian difference estimate gives
\[
 \|D_yT_x(y)\|_{\infty\to\infty}
       \le J_0CS_a^2r_V\le\tfrac14.
\]
The second-derivative bound, summed with its vertical coordinate
as root, gives
\[
 \|F(x,0)-F(0,0)\|_\infty\le CS_a\|x\|_\infty.
\]
Consequently $\|T_x(0)\|_\infty\le r_V/4$ and
\[
          \|T_x(y)\|_\infty
                   \le\tfrac14\|y\|_\infty+r_V/4\le r_V/2.
\]
Thus $T_x$ maps the complete closed ball into itself and contracts
by $1/4$. Its iterates from zero form a Cauchy sequence: consecutive
differences are bounded by a geometric sequence of ratio $1/4$.
Their limit $\rho(x)$ is a fixed point by continuity. If two fixed
points existed, their distance would be at most one quarter of
itself, so they would coincide. We have proved
\[
 \|\rho(x)\|_\infty\le r_V/2,\qquad F(x,\rho(x))=0,\qquad\rho(0)=0,
\]
with uniqueness among all stationary points in the closed
$r_V$-ball.

At each $(x,\rho(x))$ the vertical derivative of $F$ is the
invertible Hessian $H(x,\rho(x))$. The finite-dimensional smooth
implicit-function theorem therefore supplies a local smooth
stationary graph. It remains inside the open $r_V$-ball after
restricting its neighborhood, and uniqueness identifies it with
the graph just constructed. These local identifications prove
smoothness of $\rho$ throughout $\|x\|_\infty<r_X$.
If several local tags give the same function on the original
coordinate ball, the constructed graph is stationary for every
tag and is the same graph for all of them. No extension outside
the common ball is asserted.

Differentiating $F(x,\rho(x))=0$ once gives
\[
        D\rho(x)=-H(x,\rho(x))^{-1}D_xF(x,\rho(x)).
\]
The rooted row sum bounds $\|D_xF\|_{\infty\to\infty}$ by
$CS_a$. Hence $\|D\rho\|_{\infty\to\infty}\le J_0CS_a$.
Integration along a coarse line segment proves
\begin{equation}\label{eq:stationary-uniform-lipschitz}
 \|\rho(x)-\rho(x')\|_\infty
              \le C_\rho\|x-x'\|_\infty,\qquad C_\rho=1+J_0CS_a.
\end{equation}
In particular the constant used to spread fine-field slack is
independent of the auxiliary volume.

\subsection{Fixed-order derivatives with polynomial volume cost}
Let $W\ge1$ bound $|I|+|J|$. We now allow polynomial costs in $W$.
A scalar $q$-linear form with coordinate entries bounded by $C$
has supremum operator norm at most $CW^q$: expand its arguments
in coordinate basis vectors and bound the resulting at most $W^q$
terms. For a vector-valued form with supremum output norm the same
bound holds coordinate by coordinate. Passing that output to its
Euclidean norm costs at most $\sqrt W\le W$.

These elementary conversions turn
\eqref{eq:stationary-scatter-input} into polynomial bounds for the
ordinary derivatives of $f$ through order $N$. Its value has the
same kind of bound on the fixed ball, because $f(0)=0$ and the
first derivative can be integrated radially.
They also give polynomial bounds for every derivative of $\rho$
through order $N-1$. For completeness, differentiate
$F\circ\Gamma=0$, where $\Gamma(x)=(x,\rho(x))$.
At order $k\ge2$, the partition formula reads
\[
 0=H(\Gamma(x))D^k\rho(x)
   +\sum_{\substack{\pi\text{ a partition of }\{1,\ldots,k\}\\
                    |\pi|\ge2}}
       D^{|\pi|}F(\Gamma(x))
                       [D^{|B|}\Gamma(x):B\in\pi].
\]
Every block in the sum has size less than $k$; the single-block
term is exactly the isolated Hessian term, since the coarse
component of $D^k\Gamma$ vanishes. The first derivative was already
bounded in \eqref{eq:stationary-uniform-lipschitz}. Inductively,
if the lower derivatives have polynomial bounds $P_h(W)$, the sum
is bounded by
\[
 C\sum_{\substack{\pi\vdash\{1,\ldots,k\}\\|\pi|\ge2}}
       W^{|\pi|}\prod_{B\in\pi}(1+P_{|B|}(W)).
\]
Multiplying by $J_0$ bounds $D^k\rho$. This is a fixed polynomial
in $W$ at every fixed $k$, and uses derivatives of $f$ only through
order $k+1$. The Euclidean output version follows by the preceding
factor $W$. Thus the coordinate conversion does not affect the
uniform radius or the uniform first-derivative bound.

Put $M(x)=H(\Gamma(x))$. The chain rule gives polynomial bounds
for its entry derivatives through order $N-2$. A matrix with
entries bounded by $B$ has Euclidean operator norm at most $WB$,
so these are polynomial operator bounds as well. The inverse
identity gives, for each nonempty set of derivative slots $A$,
\[
 \partial_A M^{-1}
   =-M^{-1}\sum_{\varnothing\ne B\subseteq A}
                   (\partial_BM)(\partial_{A\setminus B}M^{-1}).
\]
All inverse derivatives on the right have lower order.
Together with $\|M^{-1}\|_{2\to2}\le2/m$, this proves polynomial
bounds by induction. Finally
\[
        D\log\det M[h]=\operatorname{tr}(M^{-1}DM[h]).
\]
Repeated product differentiation and
$|\operatorname{tr}A|\le W\|A\|_{2\to2}$ give polynomial bounds
for every positive derivative of $\log\det M$ through order
$N-2$. Its centered value is bounded by integrating the first
derivative. Its uncentered value is bounded too: the eigenvalues
lie between $m/2$ and a fixed upper row-sum bound, so the sum of
their logarithms has absolute value at most a fixed constant times
$W$. In particular all determinants used here are positive.

\subsection{The loop exponent and the Gaussian comparison data}
The remaining exponent is
\[
      \psi(x,y)=\ell({\cal D}_x(y))-
                              \sum_{i\in J_{\rm bond}}\log J_H(y_i),
\]
where $J_{\rm bond}$ denotes the free bonds and $J_H$ the
single-bond Haar Jacobian. The notation separates the bond set
from the earlier real-coordinate set $J$.
On the fixed combined coordinate unit ball, the actual loop
coefficient budget bounds each coordinate derivative through the
requested order by a fixed constant times $\Lambda$, summed over
at most $W$ sites. Its exponentially summable coefficient series
converges normally there, which justifies differentiation before
summation. Conversion to operator derivatives costs at most a
further power of $W$. On the half-radius Haar chart,
Section~\ref{sec:haar-chart} bounds the derivatives of $\log J_H$.
Summing over at most $W$ free bonds and composing with the fixed
Lie-coordinate maps has polynomial cost in $W$. Thus the
$\psi$ bounds have the form
\[
                       P_0(W)+P_1(W)\Lambda
\]
for fixed nonnegative polynomials $P_0,P_1$.

Apply the preceding construction at order $N=n+3$. It supplies
bounds through order $n+2$ for $\rho$, through order $n+3$ for
$f$, and through order $n+1$ for the relative log determinant.
The loop and Haar estimates are purchased through order $n+4$.
Take coarse radius $u$ and fluctuation parameter radius $v$ with
\[
 (1+C_\rho)u+vD<\min(r_H,1),\qquad u\le r_X,
 \quad c_{\rm Lie}(C_\rho u+vD)\le r_H^{\rm Haar}/2.
\]
Here $c_{\rm Lie}$ is the fixed norm of the Lie-coordinate map
and $r_H^{\rm Haar}>0$ is the Haar-chart radius. On every segment
$(x,\rho(x)+tbz)$, with $\|x\|_\infty<u$, $|b|<v$,
$0\le t\le1$ and $\|z\|_2\le D$, these inequalities verify the
needed action, coordinate and Haar domains. The stationary
Hessian is positive with inverse covariance bound $2/m$, and
stationarity removes the linear vertical Taylor term.
These are exactly the data required in
Section~\ref{sec:truncated-stationary-log}. For its final estimate
through order $j$ with the even second-order gain, choose $n\ge j+2$;
the two additional parameter derivatives must be included when fixing
the stationary-chart order.

For the selected auxiliary volumes, $W$ is bounded by a polynomial
in the auxiliary size parameter. All the prices obtained here are
therefore polynomial in that parameter, except for the single
linear loop factor. With $\Lambda\le2P_F$ their common bound can
be written $p_0+p_1P_F$, with polynomially bounded $p_0,p_1$.
This verifies the precise price dependence used in
Section~\ref{sec:stationary-exterior-receipt}. The remaining
quantitative inputs are the local-action estimate
\eqref{eq:stationary-scatter-input}, the initial Hessian gap, and
the loop coefficient and Haar-chart budgets. The construction
never changes the selected finite-volume action or its tags.
